%==============================================================================
%  CAPÍTULO VIII
%==============================================================================
\chapter{Procedimiento concursal de liquidación simplificada}
\label{cap:liquidacion-simplificada}

\fichaCapitulo{Realización rápida para personas y MIPEs.}{Entrega voluntaria,
desformalización de la incautación, venta electrónica y fin de la cuenta tradicional.}

%==============================================================================
\bloqueTeorico
%==============================================================================

\subsection{La humanización del concurso}
\pendiente{Desarrollar: desformalización de los actos de incautación en el hogar de la
persona deudora.}

\subsection{Del apremio a la colaboración: la entrega voluntaria de bienes}
\pendiente{Desarrollar: reducción de costos y sobreseimiento rápido del deudor de buena
fe.}

%==============================================================================
\bloqueNormativo
%==============================================================================

\subsection{Solicitud de liquidación voluntaria simplificada}

Exigencia de declaración jurada de antecedentes fidedignos, regulada en la \ncg{22}
\parencite{ncg-22}. Véase también \textcite{superir-faq-liquidacion-simplificada}.

\subsection{Exención de la incautación física en el hogar del deudor}

Regla general de reemplazo por entrega voluntaria y sus excepciones: orden expresa del
tribunal, infracción de los deberes del deudor como depositario provisional, o detección
de activos no declarados.

\subsection{Realización de activos mediante subasta electrónica}

Venta en plataformas electrónicas autorizadas por la \superir{}, prescindiendo del
martillero concursal tradicional.

\subsection{Supresión de la cuenta definitiva y sobreseimiento}

Exención del trámite de cuenta definitiva en procedimientos sin bienes realizables y
supresión del plazo de espera para declarar el sobreseimiento
definitivo.\verificar{Confirmar el alcance de la supresión del plazo de dos años.}

%==============================================================================
\bloquePractico
%==============================================================================

\subsection{Tratamiento de las remuneraciones del deudor persona natural}

Conforme a las instrucciones de la \superir{}, el liquidador debe ceñirse a las
siguientes directrices respecto de la remuneración del deudor
\parencite{instructivo-i5}.\verificar{Confirmar
número, fecha y contenido vigente del instructivo sobre incautación de remuneraciones.}

\begin{flujograma}{Tratamiento de la remuneración del deudor en la liquidación simplificada}{cap08-remuneraciones}
  \node[hito] (eval) {EVALUACIÓN DE LA REMUNERACIÓN DEL DEUDOR};

  \coordinate (split) at ($(eval.south)+(0,-6mm)$);
  \node[favorable, below left=12mm and 8mm of eval.south] (hasta)
    {HASTA \uf{56}\\[2pt]\normalfont Inembargable absoluto: prohibición de incautar};
  \node[adverso, below right=12mm and 8mm of eval.south] (mas)
    {MÁS DE \uf{56}\\[2pt]\normalfont Tramo excedente sujeto a incautación};

  \node[nodo, text width=34mm, below=9mm of hasta] (hastadesc)
    {Descuentos obligatorios y voluntarios: excluidos de la base};
  \node[nodo, text width=34mm, below=9mm of mas] (masdesc)
    {Obligatorios excluidos; voluntarios incluidos en la base del excedente};

  \draw[flecha] (eval.south) -- (split) -| (hasta.north);
  \draw[flecha] (eval.south) -- (split) -| (mas.north);
  \draw[flecha] (hasta) -- (hastadesc);
  \draw[flecha] (mas) -- (masdesc);
\end{flujograma}

\begin{practica}[Controversias con el empleador]
Si surge discrepancia interpretativa entre el liquidador y el empleador respecto de la
retención de los excedentes, el asunto debe someterse al tribunal del concurso por vía
incidental.\verificar{Confirmar la vía incidental aplicable y su artículo.}
\end{practica}

%==============================================================================
\bloqueJurisprudencial
%==============================================================================

\subsection{Denegación del \discharge{} por omisión de activos realizables}

Sentencias que facultan al tribunal para declarar un \discharge{} parcial o denegarlo
íntegramente cuando el deudor oculta cuentas de ahorro previsional voluntario o no
declara vehículos motorizados vigentes, por infracción de la buena fe
objetiva.\verificar{Individualizar los fallos y precisar si la denegación parcial tiene
sustento en el texto vigente o es construcción pretoriana.}
